\section{Introduction}

Human life is a sequence of decisions whose outcomes are rarely certain. These decisions run the gamut from quantitative choices, such as whether to buy a lottery ticket, to qualitative ones, such as which restaurant to try or which medical treatment to accept. Although both kinds of choices require the chooser to weigh uncertain alternatives against a reference point, the behavioural and computational literature has concentrated almost exclusively on the quantitative case. Our understanding of how people trade off the desirability of different qualitative outcomes when probabilities are imperfectly known remains correspondingly thin.

Economic and psychological theory distinguishes two components of uncertainty that jointly shape choice: risk, which arises when the probabilities of possible outcomes are precisely known, and ambiguity, which arises when these probabilities are only partially or entirely unknown \citep{ellsberg1961risk,trimmer2011bayesians}. Individuals generally exhibit aversion to both risk and ambiguity when choosing between gains, but these two aversions vary substantially across individuals and are only weakly correlated with each other \citep{levy2010neural,tymula2013lifespan}. Characterising these attitudes at the individual level is important because they predict real-world behaviour and clinical phenotypes: ambiguity-related choice anomalies have been linked with aging \citep{tymula2013lifespan}, post-traumatic stress disorder in combat veterans \citep{ruderman2016ptsd}, and clinical decision-making such as breast cancer screening uptake \citep{yeh2015dense} and prostate cancer screening decisions \citep{james2017prostate}.

The dominant experimental paradigm for measuring risk and ambiguity attitudes presents participants with choices between a small certain gain and a lottery whose monetary payoff and probability of winning are parametrically varied \citep{levy2010neural,levy2012measuring,jia2020learning}. The analytic backbone is a classical utility function in which a single risk-attitude parameter governs the curvature of the mapping from dollars to subjective value. This machinery fits monetary choices well but presupposes that every outcome can be placed on a single numeric axis. When the choice set is instead qualitative---for example, the severity of symptoms after a medical treatment, or the type of partner one spends an evening with---the utility function can no longer be fit directly. The typical workarounds are either to quantify some surrogate feature of the outcome (mg of medication, months of extended lifespan, millilitres of water) or to treat qualitative categories as equally spaced \citep{bijlenga2007balance,basu2018vio}. Both strategies abandon the original qualitative framing of the choice and discard information about individual differences in how people weight each category.

This gap matters for clinical, neuroimaging and computational-psychiatry research. Medical diagnoses and treatment outcomes are intrinsically qualitative \citep{basu2018vio,james2017prostate}, yet existing computational tools force them into one of the two surrogates above. Similarly, investigations that link choice behaviour to anxiety, depression, obsessive-compulsive disorder or binge eating disorder \citep{ruderman2016ptsd,voon2015bed} would benefit from a framework that accommodates qualitative outcomes directly, rather than shoehorning them into monetary equivalents.

In this work we develop such a framework. We ran two samples---an in-person cohort of 66 cognitively healthy adults screened with the Montreal Cognitive Assessment \citep{nasreddine2005moca,fujiwara2010mocaj} and an online replication cohort of 332 participants---through a decision-making task in which the same individuals faced both monetary choices and hypothetical medical choices calibrated to the same probability and ambiguity levels \citep{levy2012measuring,jia2020learning}. We fit hierarchical Bayesian models using the No-U-Turn Sampler implemented in PyMC \citep{hoffman2014nuts,salvatier2016pymc3} and compared models using leave-one-out cross-validation with ArviZ \citep{kumar2019arviz,vehtari2017loo}. Our Estimated Value model replaces the single utility-curve parameter with an individualised set of subjective values, one for each ordered qualitative outcome, estimated directly from the choices. Our contributions are:
\begin{itemize}
  \item An Estimated Value hierarchical Bayesian model that extracts individualised subjective values for ordinal qualitative outcomes without assuming equal spacing.
  \item Direct validation: in the monetary domain, where objective values are available, the Estimated Value model outperforms a classical utility function and its trembling-hand extension in both samples.
  \item Evidence that ambiguity aversion ($\beta$) generalises across domains: medical $\beta$ is strongly associated with monetary $\beta$ in-person (mean slope = 0.76, 89\% HDPi [0.63, 0.91]) and online (mean slope = 0.37, 89\% HDPi [0.29, 0.44]).
  \item Simulations that distinguish genuine model advantage from overfitting by varying the noise on the latent generative utility function.
\end{itemize}

\section{Related Work}

\paragraph{Risk and ambiguity attitudes under monetary outcomes.}
A now-standard line of work has used parametric gambling tasks to measure risk and ambiguity attitudes in healthy adults and clinical populations. \citet{levy2010neural} established that subjective value under risk and under ambiguity share a common neural representation in striatum and medial prefrontal cortex while the two attitudes are behaviourally dissociable. \citet{levy2012measuring} formalised the accompanying experimental-economics protocol---participants choose between a certain gain and a lottery that varies in either outcome probability or ambiguity level, and a power utility function is fit to their choices. This protocol has since been extended to trace lifespan changes \citep{tymula2013lifespan}, the effects of learning interventions \citep{jia2020learning}, sleep deprivation \citep{mullette2015sleep}, and differences induced by brain lesions \citep{clark2008vmpfc}. Across this literature, attitudes towards risk and ambiguity are treated as stable individual-difference measures, though their test-retest and cross-domain stability has been debated \citep{jia2020learning}.

\paragraph{Qualitative and medical outcomes.}
Attempts to characterise uncertainty attitudes outside the monetary domain have relied on numeric surrogates. Studies of medical decision-making quantify outcomes as months of extended lifespan \citep{basu2018vio}, values placed on health states at term pregnancy \citep{bijlenga2007balance}, or willingness to undergo screening \citep{yeh2015dense,james2017prostate}. While these approaches are useful for specific clinical questions, they collapse the qualitative structure of the outcome space onto a single axis. Consequently, they cannot capture how a person weights, for example, ``moderate improvement'' relative to ``complete recovery'' on that person's own scale.

\paragraph{Clinical relevance of ambiguity attitudes.}
Ambiguity aversion has been linked to clinical and subclinical outcomes. \citet{ruderman2016ptsd} showed that combat veterans with PTSD exhibit heightened aversion specifically to ambiguous losses, and that the degree of aversion mediates the link between combat exposure and anxious arousal. \citet{voon2015bed} reviews how decision-making under ambiguity is reshaped across obesity and binge eating disorder. Dense-breast notifications reduce screening intentions more strongly in ambiguity-averse women \citep{yeh2015dense}, and uncertainty about the value of the procedure shapes prostate-cancer screening decisions \citep{james2017prostate}. A framework that extracts individualised subjective values from qualitative medical choices would let this literature move beyond surrogate quantifications.

\paragraph{Hierarchical Bayesian modelling of choice.}
Hierarchical Bayesian modelling has become a standard tool for estimating individual-level parameters from behavioural choice data because it exploits partial pooling across participants to stabilise individual estimates \citep{salvatier2016pymc3,kumar2019arviz}. Posterior inference is routinely performed using the No-U-Turn Sampler \citep{hoffman2014nuts}, and competing models are compared with leave-one-out cross-validation on the expected log pointwise predictive density \citep{vehtari2017loo}. We build directly on this tooling and on the power-utility formulation of \citet{levy2012measuring}. Our contribution is to replace the single-parameter utility curve with an individualised vector of estimated subjective values that can be fit equally well to quantitative and qualitative outcomes, and to show that doing so does not cost predictive accuracy---in fact it improves it, even when objective dollar amounts are available.

\paragraph{Foundational theories of value.}
Finally, our framework connects to two foundational ideas in the psychology of value. \citet{kahneman1979prospect} showed that the mapping from objective outcomes to subjective value is non-linear and frame-dependent, and \citet{parducci1965range} demonstrated that category judgements depend on the range and frequency of outcomes encountered. A single-parameter utility curve cannot easily accommodate either phenomenon. Our Estimated Value model, by assigning a separate posterior to each outcome category, provides a structure in which the context dependence predicted by range-frequency theory and the curvature differences predicted by prospect theory can emerge naturally from the data.
